% ─────────────────────────────────────────────────────────────────────
%  Capítulo 6 — Resultados y discusión
% ─────────────────────────────────────────────────────────────────────
\chapter{Resultados y discusión}
\label{cap:resultados}

Este capítulo recoge los resultados experimentales y, sobre todo, los discute. El recorrido es el mismo que el del Cap.~\ref{cap:diseno}: empezamos por el rendimiento individual de cada detector, miramos el efecto del postprocesado temporal, desglosamos por tipo y severidad de ataque, hablamos de la recalibración de umbrales, comentamos lo que aporta el \textit{ensemble} y lo comparamos con el techo supervisado. Cierran el capítulo los análisis transversales: explicabilidad SHAP, generalización a UK-DALE, robustez frente a atacantes adaptativos y viabilidad de despliegue.

Todas las cifras salen de los ficheros \code{data/metrics\_*.json} regenerados al ejecutar cada notebook con \code{nbconvert --execute}, así que son reproducibles bit a bit fijando las semillas declaradas.

\section{Detectores individuales: tabla maestra}
\label{sec:resultados-individuales}

\begin{table}[H]
\centering
\caption{Tabla maestra consolidada: rendimiento y coste operativo de cada detector sobre \code{test\_with\_attacks} con los hiperparámetros calibrados en validación. Métricas estadísticas \emph{smoothed} a nivel de minuto; latencia mediana de detección por episodio; tamaño en disco del modelo serializado; ms/muestra medidos en Apple M2~Pro (CPU, batch=1).}
\label{tab:master}
\scriptsize
\begin{tabularx}{\textwidth}{X r r r r r r r r}
\toprule
Detector & $\fone$ & $\ftwo$ & Prec. & Rec. & $\aucpr$ & Lat.\,(min) & Tam.\,(kB) & ms/m.\\
\midrule
IF                 & 0{,}754 & 0{,}748 & 0{,}762 & 0{,}747 & 0{,}782 & 2{,}0 &  3\,493 & 0{,}005 \\
K-Means+IF         & 0{,}790 & 0{,}768 & 0{,}834 & 0{,}750 & 0{,}764 & 4{,}0 & 17\,425 & 0{,}008 \\
\textbf{Dense-AE}  & \textbf{0{,}816} & 0{,}809 & 0{,}827 & 0{,}805 & 0{,}855 & 2{,}0 & \textbf{134} & \textbf{0{,}001} \\
LSTM-AE            & 0{,}745 & 0{,}789 & 0{,}671 & 0{,}838 & 0{,}791 & 0{,}0 &    737 & 0{,}003 \\
CNN-AE             & 0{,}435 & 0{,}449 & 0{,}366 & 0{,}537 & 0{,}432 & 0{,}0 &    199 & 0{,}001 \\
Transformer-AE     & 0{,}604 & 0{,}633 & 0{,}530 & 0{,}700 & 0{,}584 & 0{,}0 &  1\,103 & 0{,}001 \\
LightGBM$^\star$   & 0{,}910 & 0{,}871 & 0{,}954 & 0{,}871 & 0{,}941 & 0{,}0 &    125 & 0{,}0002 \\
Stacking           & 0{,}811 & 0{,}803 & 0{,}820 & 0{,}803 & 0{,}881 & 0{,}0 &  4\,795 & 0{,}012 \\
\bottomrule
\multicolumn{9}{l}{\footnotesize $^\star$ Techo supervisado de referencia, no detector final propuesto. Negrita: detector recomendado para despliegue \textit{edge}.}
\end{tabularx}
\end{table}

\paragraph{Lectura de la tabla.} De la \cref{tab:master} salen varias lecturas que conviene comentar.

Empezando por arriba, \textbf{LightGBM domina}, como cabía esperar: al ver etiquetas durante el entrenamiento alcanza $\aucpr=0{,}941$, claramente por encima de los detectores no supervisados. Entre estos, \textbf{el Dense-AE es el ganador} tanto en $\fone$ (0{,}816) como en $\aucpr$ (0{,}855), con un equilibrio excelente entre precisión (0{,}827) y \textit{recall} (0{,}805); el \textit{bottleneck} agresivo de 17$\to$3 parece ser, en este dataset, un \textit{sweet spot}. El \textbf{LSTM-AE} se queda por detrás, con $\aucpr=0{,}791$ y $\fone=0{,}745$; consigue un \textit{recall} muy alto (0{,}838) pero su precisión cae (0{,}671) porque la calibración prioriza $F_2$, algo que la recalibración por $F_1$ del \cref{sec:resultados-recalibracion} corrige parcialmente. En este experimento, la arquitectura recurrente añade coste sin aportar una ventaja clara sobre el AE denso.

Los modelos clásicos (IF y K-Means+IF) saturan en torno a $\aucpr\approx 0{,}78$, lo que evidencia el techo de un detector que no aprende representación. El \textbf{Transformer-AE} queda por debajo del LSTM-AE, probablemente porque con sólo 78\,k parámetros y dos bloques no aprovecha el sesgo inductivo de la atención; troceando ventanas de 60 minutos en \textit{patches} de 5, se obtienen únicamente 12 \textit{tokens}, demasiado pocos para que la atención marque diferencia. El \textbf{CNN-AE} es el peor del banco: incluso después de corregir el bug de los pesos (\cref{sec:nb07}), que llevó el $F_1$ de 0{,}27 a 0{,}43, sigue lejos del resto; una arquitectura sólo convolucional, sin contexto recurrente ni atención, no parece dar para más en este problema.

Finalmente, el \textbf{Stacking} llega a $\fone=0{,}811$, lo que lo coloca por encima de varios detectores individuales pero ---inesperadamente--- ligeramente por debajo del Dense-AE. Sobre por qué ocurre esto volvemos en el \cref{sec:resultados-ensemble}.

Las curvas ROC y Precision-Recall completas de los ocho detectores aparecen superpuestas en la \cref{fig:roc-pr}, donde se aprecia visualmente la dominancia de LightGBM y el orden relativo entre los detectores no supervisados.

\begin{figure}[H]
  \centering
  \includegraphics[width=0.95\textwidth]{roc_pr_all.png}
  \caption{Curvas ROC y Precision-Recall de los ocho detectores sobre \code{test\_with\_attacks}.}
  \label{fig:roc-pr}
\end{figure}

Para facilitar la comparación visual de los cuatro mejores detectores (LightGBM, Dense-AE, LSTM-AE y el \emph{Stacking}) sobre las cuatro métricas principales, la \cref{fig:comparison} agrupa los valores en un único diagrama de barras. Esta forma de visualización condensa la tabla maestra en un golpe de vista y revela los compromisos: el Dense-AE iguala la precisión del LightGBM pero a costa de un \textit{recall} ligeramente menor; el LSTM-AE compensa una precisión modesta con el mejor \textit{recall} entre los no supervisados; el \emph{Stacking} queda en una posición intermedia sin destacar en ninguna métrica.

\begin{figure}[H]
  \centering
  \includegraphics[width=0.92\textwidth]{comparison_4_models.png}
  \caption{Comparación de los cuatro detectores destacados sobre las cuatro métricas principales ($F_1$, precision, recall, $\aucpr$). El LightGBM como techo supervisado queda en negro; los no supervisados, en color.}
  \label{fig:comparison}
\end{figure}

\section{Efecto del postprocesado temporal}
\label{sec:resultados-temporal}

El filtro $W/K$ es simple ---una predicción se confirma sólo si hay al menos $K$ positivos en los últimos $W$ minutos--- pero, como muestra la \cref{tab:wk-effect}, su efecto sobre cada detector es muy distinto:

\begin{table}[H]
\centering
\caption{Efecto del postprocesado temporal: $F_1$ raw frente a $F_1$ smoothed con el $(W, K)$ óptimo en validación.}
\label{tab:wk-effect}
\footnotesize
\begin{tabularx}{\textwidth}{l r r r r X}
\toprule
Detector & $F_1$ raw & $F_1$ smoothed & $\Delta F_1$ & $(W, K)$ & Comentario \\
\midrule
IF                 & 0{,}000 & 0{,}754 & +0{,}754 & (3, 2)   & Recupera detección incluso con thr extremo \\
K-Means+IF         & 0{,}000 & 0{,}790 & +0{,}790 & (11, 10) & Consenso casi unánime en 11 min \\
Dense-AE           & 0{,}781 & 0{,}816 & +0{,}035 & (3, 3)   & Smoothing fino \\
LSTM-AE            & 0{,}744 & 0{,}745 & +0{,}001 & (3, 3)   & Smoothing apenas mejora — ver discusión \\
CNN-AE             & 0{,}431 & 0{,}435 & +0{,}004 & (15, 15) & Consenso amplio insuficiente por baja AUC \\
Transformer-AE     & 0{,}596 & 0{,}604 & +0{,}008 & (11, 11) & Mejora marginal \\
\bottomrule
\end{tabularx}
\end{table}

La cuantía de la mejora depende sobre todo de cómo de bueno fuera el \textit{raw}. En el IF y el K-Means+IF, las predicciones \textit{raw} son ráfagas de positivos sueltos y el filtro mejora muchísimo la precisión. En los AE profundos, en cambio, el \textit{raw} ya es bastante competitivo y el postprocesado apenas añade nada.

\paragraph{Visualización temporal del efecto.} La \cref{fig:timelines} muestra, en formato \emph{timeline}, el comportamiento de cuatro detectores representativos (IF, K-Means+IF, Dense-AE y LSTM-AE) sobre un fragmento del test que contiene varios episodios de ataque. Sobre cada \emph{timeline} se superponen los segmentos verdaderos de ataque (bandas rojas) y los disparos del detector (líneas verticales). El IF, sin postprocesado, produce un ruido casi continuo de alarmas puntuales; tras el filtro $W/K$, esa actividad se concentra en los episodios reales. Los AE muestran patrones más limpios desde el principio.

\begin{figure}[H]
  \centering
  \begin{subfigure}[t]{0.48\textwidth}
    \centering
    \includegraphics[width=\linewidth]{timeline_if.png}
    \caption{Isolation Forest.}
  \end{subfigure}\hfill
  \begin{subfigure}[t]{0.48\textwidth}
    \centering
    \includegraphics[width=\linewidth]{timeline_k-means_if.png}
    \caption{K-Means + IF.}
  \end{subfigure}\\[1mm]
  \begin{subfigure}[t]{0.48\textwidth}
    \centering
    \includegraphics[width=\linewidth]{timeline_dense_ae.png}
    \caption{Dense Autoencoder.}
  \end{subfigure}\hfill
  \begin{subfigure}[t]{0.48\textwidth}
    \centering
    \includegraphics[width=\linewidth]{timeline_lstm_ae.png}
    \caption{LSTM Autoencoder.}
  \end{subfigure}
  \caption{Comportamiento temporal de cuatro detectores sobre el conjunto de test. Bandas rojas: episodios de ataque verdaderos; líneas verticales: alarmas del detector tras postprocesado $W/K$. Los detectores clásicos requieren un postprocesado más agresivo para limpiar el ruido de fondo; los AE producen detecciones más concentradas en los episodios reales.}
  \label{fig:timelines}
\end{figure}

\paragraph{Detalle a una semana.} Un \emph{zoom} sobre una sola semana (\cref{fig:timelines-zoom}) permite ver cómo cada detector se comporta frente a la variabilidad natural del consumo doméstico (fines de semana, ciclos diarios) y dónde sitúa exactamente sus alarmas respecto al inicio del ataque.

\begin{figure}[H]
  \centering
  \begin{subfigure}[t]{0.48\textwidth}
    \centering
    \includegraphics[width=\linewidth]{timeline_zoom_if.png}
    \caption{IF (zoom semanal).}
  \end{subfigure}\hfill
  \begin{subfigure}[t]{0.48\textwidth}
    \centering
    \includegraphics[width=\linewidth]{timeline_zoom_dense_ae.png}
    \caption{Dense-AE (zoom semanal).}
  \end{subfigure}
  \caption{\emph{Zoom} de una semana del test sobre dos detectores. La concentración temporal de las detecciones es notablemente distinta: el IF dispara con persistencia constante; el Dense-AE responde con \emph{ráfagas} más discriminantes alineadas con los episodios.}
  \label{fig:timelines-zoom}
\end{figure}

\section{Análisis por tipo y severidad de ataque}
\label{sec:resultados-tipo-severidad}

Esta es, posiblemente, la sección que más información aporta sobre el comportamiento real de cada detector. La \cref{fig:heatmap-recall} desglosa el \textit{recall} por tipo de ataque y severidad para los dos mejores detectores no supervisados (Dense-AE y LSTM-AE) y para el techo supervisado (LightGBM):

\begin{figure}[H]
  \centering
  \includegraphics[width=0.95\textwidth]{heatmap_recall.png}
  \caption{Recall por tipo de ataque y severidad. Filas: tipos. Columnas: \emph{low}, \emph{medium}, \emph{high}.}
  \label{fig:heatmap-recall}
\end{figure}

\paragraph{Lo que se observa.} La severidad \emph{high} es un caso fácil casi para cualquier detector y tipo: el \textit{recall} se mantiene cerca de 1. En la severidad \emph{low}, en cambio, los números bajan típicamente al rango 0{,}50--0{,}80, lo que confirma que los ataques sutiles entran parcialmente en la variabilidad natural del consumo.

Yendo a casos concretos, el \code{voltage\_spoof} es detectado prácticamente al 100\,\% por Dense-AE y LightGBM en las tres severidades (recall $\geq$0{,}98), lo que confirma directamente el valor diagnóstico de $\vires$. El \code{noise} es donde brilla el LSTM-AE (\textit{recall} 1{,}00 en todas las severidades): la red recurrente capta la incoherencia local en la serie con una sensibilidad que las otras arquitecturas no igualan. Y luego está el \code{step}, el caso más difícil del banco, sobre todo en severidad \emph{low} y \emph{medium}: un salto brusco aislado se reabsorbe rápido en las características de ventana corta, hasta el punto de que ni siquiera LightGBM (\textit{recall} 0{,}60--0{,}65) consigue superar a los AE. Los ataques \code{replay} y \code{ramp} quedan en dificultad intermedia; el \code{replay} merece mención aparte, porque ahí el LSTM-AE sí saca ventaja real (\textit{recall} 0{,}77/0{,}90/0{,}99 frente a Dense-AE 0{,}57/0{,}77/0{,}87) gracias a su modelización temporal explícita.

Hay dos lecturas que sacar de aquí. Una, metodológica: distintos detectores son fuertes en categorías distintas, lo que justifica el enfoque de \textit{ensemble}. Otra, operativa: en producción no basta con un único umbral global, conviene reportar el \textit{recall} esperado por tipo y severidad para que el operador sepa de qué se fía y de qué no.

\section{Recalibración de umbrales}
\label{sec:resultados-recalibracion}

El NB11 recorre el espacio de calibración variando dos ejes: la métrica objetivo ($F_1$, $F_2$, $F_{0.5}$ o un coste FN/FP explícito) y el tipo de búsqueda (umbral fijo o \textit{joint search} con $W$ y $K$). El caso más ilustrativo es el del LSTM-AE: su calibración por defecto va a por $F_2$, lo que prioriza \textit{recall} (0{,}838) a costa de precisión (0{,}671). Una recalibración con \textit{joint search} maximizando $F_1$ \textit{smoothed} consigue subir el $F_1$ de 0{,}65 a 0{,}75, con precisión 0{,}65 y \textit{recall} 0{,}86.

\paragraph{Recomendación operativa.} El umbral óptimo no es uno solo: depende del coste relativo de FP y FN en el escenario de despliegue. En detección de fraude, un FN deja escapar dinero y un FP genera una investigación inútil; la razón típica en operación industrial ronda 5:1 ---un FN cuesta cinco veces más que un FP---, lo que sugiere optimizar $F_{1{,}5}$ o $F_2$. En cambio, ante un escenario de ataque coordinado las consecuencias de un FN son mucho peores y conviene irse a $F_3$ o, directamente, a un coste asimétrico explícito.

\section{Stacking}
\label{sec:resultados-ensemble}

El \emph{stacking} con Meta-LightGBM produce el siguiente resumen interno:

\begin{table}[H]
\centering
\caption{Estrategias de combinación evaluadas en NB10. Métricas sobre validación, antes de la selección final.}
\label{tab:stacking}
\small
\begin{tabularx}{\textwidth}{l r r r}
\toprule
Estrategia & $\aucroc$ & $\aucpr$ & $\fone$ \\
\midrule
Mean blending     & 0{,}890 & 0{,}594 & 0{,}635 \\
Weighted blending & 0{,}927 & 0{,}764 & 0{,}732 \\
Meta-LR           & 0{,}934 & 0{,}792 & 0{,}756 \\
Meta-LightGBM     & 0{,}971 & 0{,}920 & 0{,}842 \\
\bottomrule
\end{tabularx}
\end{table}

\paragraph{Discusión crítica.} El Meta-LightGBM se impone con claridad a las alternativas lineales, ahí no hay sorpresa. Lo interesante es que, ya en test, el \textit{stacking} cae a $\fone=0{,}811$, ligeramente por debajo del Dense-AE individual ($\fone=0{,}816$). ¿Por qué?

Las hipótesis no se excluyen entre sí. La más probable es que los detectores base están correlados en sus errores: todos fallan en los mismos ataques difíciles ---típicamente \textit{replay} en severidad \textit{low}---, así que la diversidad efectiva del \textit{ensemble} es menor de lo que sugería el diseño. También cabe que el meta-learner sobreajuste a validación: con $\sim$144\,k muestras y muchos parámetros, hay margen para que aprenda particularidades de \textit{val} que no se trasladan a \textit{test}. Y, en última instancia, puede que las anomalías genuinamente difíciles requieran información que ninguno de los detectores base captura ---contexto externo como ocupación, tarifa horaria o calendario laboral, por ejemplo.

\paragraph{Conclusión metodológica.} Un \textit{stacking} no siempre mejora al mejor detector individual, y conviene reportarlo así. Mostrar el resultado tal cual y discutir el porqué es mucho mejor que presentar un \textit{ensemble} marginalmente mejor sin contextualizar lo que está pasando.

\section{Comparación con el techo supervisado}
\label{sec:resultados-techo}

LightGBM llega a $\fone=0{,}910$ y $\aucpr=0{,}941$ frente a los $\fone=0{,}816$ y $\aucpr=0{,}855$ del mejor no supervisado (Dense-AE). La brecha ronda los 9 puntos en ambas métricas.

Esa brecha no es un número suelto: cuantifica el coste de renunciar a las etiquetas. Con un corpus de ataques reales etiquetados, parte de esa diferencia se podría cerrar; pero ese corpus no existe en producción, y los detectores no supervisados se quedan en torno al 90\,\% del techo, lo que es operativamente más que aceptable. Esta es, en el fondo, la tesis que defiende todo el TFG: los detectores no supervisados están cerca del óptimo alcanzable sin etiquetas, y la inversión en etiquetar ataques ---con su coste ético y económico--- ofrecería un retorno limitado.

\section{Explicabilidad SHAP}
\label{sec:resultados-shap}

El análisis SHAP sobre LightGBM revela las características más influyentes en la decisión:

\begin{table}[H]
\centering
\caption{Top-10 características por importancia SHAP (\textit{TreeExplainer} exacto sobre LightGBM, muestra de 5\,000 minutos de \code{test\_with\_attacks}). $|\bar\phi_j|$ en unidades del \emph{log-odds} de salida; valores normalizados a la suma del Top-10 entre paréntesis.}
\label{tab:shap-top10}
\small
\begin{tabularx}{0.75\textwidth}{X r r}
\toprule
Característica & $|\bar\phi_j|$ & \% Top-10 \\
\midrule
\code{vi\_res\_roll15\_mean}    & 0{,}7929 & 43{,}3\,\% \\
\code{gap\_intensity\_ratio}    & 0{,}2925 & 16{,}0\,\% \\
\code{vi\_res\_abs}             & 0{,}1424 &  7{,}8\,\% \\
\code{hour\_cos}                & 0{,}1259 &  6{,}9\,\% \\
\code{sm\_gap\_ratio}           & 0{,}0830 &  4{,}5\,\% \\
\code{Global\_active\_power}    & 0{,}0655 &  3{,}6\,\% \\
\code{Global\_reactive\_power}  & 0{,}0594 &  3{,}3\,\% \\
\code{Global\_intensity}        & 0{,}0532 &  2{,}9\,\% \\
\code{dow\_sin}                 & 0{,}0513 &  2{,}8\,\% \\
\code{hour\_sin}                & 0{,}0416 &  2{,}3\,\% \\
\bottomrule
\end{tabularx}
\end{table}

\paragraph{Cuando SHAP coincide con la física.} Las tres características que SHAP coloca arriba son precisamente las que cualquier ingeniero familiarizado con el problema señalaría: la media móvil del residual $\vires$ (que se lleva el 43\,\% de la importancia del Top-10), el cociente potencia/intensidad y la magnitud absoluta de $\vires$. Esa coincidencia entre lo que el modelo aprende y lo que la intuición física dice que debería importar funciona como una validación cruzada: el modelo no está pillando correlaciones espurias sino la estructura física del problema. Y explica también por qué la dominancia tan acusada de \code{vi\_res\_roll15\_mean} hace detectable cualquier ataque que rompa $P\approx V\cdot I$, y por qué el mismo modelo entrenado sin esa característica (modo \code{base}) pierde varios puntos de $\fone$. La \cref{fig:shap-beeswarm} muestra la distribución completa de los valores SHAP por característica: el ancho de cada banda indica la dispersión de la atribución, y el color, el valor de la característica.

\begin{figure}[H]
  \centering
  \begin{subfigure}[t]{0.48\textwidth}
    \centering
    \includegraphics[width=\linewidth]{shap_summary_bar_lightgbm.png}
    \caption{Importancia media absoluta por característica.}
    \label{fig:shap-bar}
  \end{subfigure}\hfill
  \begin{subfigure}[t]{0.48\textwidth}
    \centering
    \includegraphics[width=\linewidth]{shap_beeswarm.png}
    \caption{Distribución (\emph{beeswarm}) de los valores SHAP.}
    \label{fig:shap-beeswarm}
  \end{subfigure}
  \caption{Atribución SHAP sobre el modelo LightGBM (NB12). El \emph{bar plot} (a) ordena las características por importancia media; el \emph{beeswarm} (b) revela además la dispersión y la dirección del impacto: el color codifica el valor de la característica.}
  \label{fig:shap-global}
\end{figure}

\paragraph{Comparación con SHAP sobre Dense-AE (no supervisado).} Aplicar el mismo análisis sobre el mejor detector no supervisado (Dense-AE) permite contrastar qué características pesan en un modelo que \emph{nunca ha visto etiquetas de ataque}. La \cref{fig:shap-dense} muestra el resultado: las características dominantes coinciden cualitativamente con las de LightGBM (los derivados de $\vires$ y los cocientes físicos), pero el orden y la magnitud relativa cambian. Esta convergencia entre un modelo supervisado y otro completamente no supervisado refuerza la robustez de la conclusión: el aprendizaje, con o sin etiquetas, identifica la misma estructura física como discriminante.

\begin{figure}[H]
  \centering
  \includegraphics[width=0.85\textwidth]{shap_summary_dense_ae.png}
  \caption{Importancia SHAP sobre el Dense-AE no supervisado. Comparar con la \cref{fig:shap-bar}: las características dominantes son cualitativamente las mismas pese a la ausencia total de etiquetas durante el entrenamiento.}
  \label{fig:shap-dense}
\end{figure}

\paragraph{Análisis por tipo de ataque.} El \emph{heatmap} SHAP $\times$ tipo de ataque (\cref{fig:shap-per-attack}) confirma que \code{vi\_res\_roll15\_mean} y \code{gap\_intensity\_ratio} son las dos características dominantes para los siete tipos, pero el tercer puesto varía y resulta diagnóstico: en \code{voltage\_spoof} entra \code{Voltage} en el podio (única familia que ataca directamente la tensión); en \code{offset}, \code{ramp} y \code{step} entra \code{Global\_active\_power} (ataques aditivos sobre la potencia); y en \code{scaling}, \code{noise} y \code{replay} se mantiene \code{vi\_res\_abs} (ataques que rompen la invariante por escalado o ruido). Esta complementariedad es consistente con la intuición física y refuerza la coherencia del modelo.

\begin{figure}[H]
  \centering
  \includegraphics[width=0.95\textwidth]{shap_per_attack_type.png}
  \caption{\emph{Heatmap} de importancia SHAP por (tipo de ataque, característica) sobre LightGBM. Filas: tipos de ataque; columnas: características. Intensidad proporcional a \texttt{mean$|$SHAP$|$}.}
  \label{fig:shap-per-attack}
\end{figure}

\paragraph{Explicaciones locales.} Más allá del análisis global, SHAP permite explicar \emph{cada predicción individual}. La \cref{fig:shap-locales} muestra tres casos representativos: un acierto (TP, ataque correctamente detectado), un falso positivo (FP, normal etiquetado como ataque) y un falso negativo (FN, ataque que el modelo no captura). En el TP, las características $\vires$ y derivadas empujan claramente la predicción hacia ``ataque''; en el FP, se observa una coincidencia espuria de varias características en valores extremos pero legítimos (consumo inusual pero normal); en el FN, las características clave permanecen en su rango habitual, lo que delata un ataque \emph{stealth} de baja severidad. Estas explicaciones locales son la herramienta que, en producción, permitiría a un operador humano validar o desestimar una alarma del sistema sin necesidad de reconstruir el modelo desde cero.

\begin{figure}[H]
  \centering
  \begin{subfigure}[t]{0.32\textwidth}
    \centering
    \includegraphics[width=\linewidth]{shap_decision_TP_1.png}
    \caption{Verdadero positivo (TP).}
  \end{subfigure}\hfill
  \begin{subfigure}[t]{0.32\textwidth}
    \centering
    \includegraphics[width=\linewidth]{shap_decision_FP_1.png}
    \caption{Falso positivo (FP).}
  \end{subfigure}\hfill
  \begin{subfigure}[t]{0.32\textwidth}
    \centering
    \includegraphics[width=\linewidth]{shap_decision_FN_1.png}
    \caption{Falso negativo (FN).}
  \end{subfigure}
  \caption{Explicaciones locales SHAP (gráficos de decisión) sobre tres muestras del test: un acierto, un falso positivo y un falso negativo. Cada línea es la trayectoria de la predicción acumulando las contribuciones de las características, ordenadas por magnitud.}
  \label{fig:shap-locales}
\end{figure}

\section{Generalización cross-domain (UK-DALE)}
\label{sec:resultados-crossdomain}

El NB13 mide cuánto del Dense-AE entrenado sobre UCI sigue funcionando cuando lo aplicamos a UK-DALE \cite{kelly2015ukdale}, un conjunto distinto que recoge cinco viviendas británicas con una frecuencia de muestreo mayor ---6 segundos frente a 1 minuto--- y que, además, no incluye intensidad como variable independiente, lo que obliga a estimar $\vires$ de forma aproximada.

\paragraph{Resultado.} La caída se cuantifica sin maquillaje: $\aucpr$ baja de 0{,}855 a $\sim$0{,}60 sin tocar nada. Recalibrando únicamente el umbral sobre una pequeña muestra de UK-DALE ---y manteniendo el modelo congelado--- $\aucpr$ vuelve a subir hasta $\sim$0{,}73. De aquí salen tres lecturas. La primera, que el modelo transfiere parcialmente: las regularidades del consumo doméstico son lo bastante parecidas entre los dos dominios como para que la transferencia tenga sentido. La segunda, que basta una recalibración de umbral para recuperar gran parte del rendimiento, lo que sugiere que el modelo aprende a discriminar pero no a normalizar la escala absoluta del residual. Y la tercera, ya operativa: en producción, una recalibración por dominio sería el procedimiento estándar de despliegue.

\section{Robustez frente a atacantes adaptativos}
\label{sec:resultados-adversarial}

El NB14 se hace la pregunta incómoda: ¿qué pasa si el atacante conoce el detector y se calibra justo a la frontera de detección? Hablamos por tanto de un atacante adaptativo de \textit{caja blanca}, con acceso al modelo y a los umbrales. Probamos dos escenarios. En el primero, evasión por reducción de severidad, el atacante baja el factor $\alpha$ del ataque \textit{scaling} hasta que el detector deja de dispararse; la frontera de detección típica queda en $\alpha\in[0{,}85, 0{,}90]$ (reducciones del 10--15\,\%) para los detectores individuales y se desplaza a $\alpha\in[0{,}90, 0{,}92]$ para el \textit{ensemble}. En el segundo, evasión por suavización temporal, el atacante reparte el ataque en pequeños incrementos para colarse por debajo del filtro $W/K$; esto le obliga a sostener la manipulación durante mucho más tiempo, lo que la hace operativamente menos práctica.

\paragraph{Lectura económica.} La frontera de detección no es sólo un dato técnico: define el umbral económico del fraude. Por debajo de ese punto el ataque queda invisible para el detector, pero al mismo tiempo se vuelve marginal en rentabilidad. La detección no necesita cero falsos negativos para ser útil; basta con que el coste de la evasión supere el beneficio que el atacante puede sacar.

\section{Viabilidad de despliegue \textit{edge}}
\label{sec:resultados-edge}

Las tres columnas finales de la \cref{tab:master} ---latencia, tamaño y ms/muestra--- son las que determinan la viabilidad de cada detector en un contador o concentrador inteligente. La \cref{tab:edge-clases} las traduce a clases de hardware:

\begin{table}[H]
\centering
\caption{Clasificación de detectores por viabilidad en plataformas \textit{edge} representativas. Clases: MCU = microcontrolador ARM Cortex-M con $\sim$256\,kB Flash; SBC = computador monoplaca tipo Raspberry Pi.}
\label{tab:edge-clases}
\small
\begin{tabularx}{\textwidth}{X l X}
\toprule
Detector & Plataforma viable & Comentario \\
\midrule
\textbf{Dense-AE}      & \textbf{MCU} & 134\,kB, $<$0{,}001\,ms/m, $\fone$=0{,}816 — \textbf{recomendado}. \\
LightGBM$^\star$       & MCU         & 125\,kB, requiere etiquetas (no detector final). \\
CNN-AE                 & MCU         & 199\,kB pero $\fone$ insuficiente. \\
LSTM-AE                & SBC ligero  & 737\,kB; alternativa si se prioriza $\aucpr$. \\
Transformer-AE         & SBC ligero  & 1\,103\,kB; coste no justificado en este experimento. \\
IF / K-Means+IF        & SBC potente & 3{,}5--17\,MB; sólo viable en concentrador o gateway. \\
Stacking               & SBC potente & 4{,}8\,MB; requiere ejecutar los detectores base en paralelo. \\
\bottomrule
\end{tabularx}
\end{table}

\paragraph{Conclusión.} El Dense-AE es nuestra recomendación para despliegue en contador: 134\,kB en disco, menos de 1\,$\mu$s por muestra de inferencia en CPU M2 Pro y $\fone=0{,}816$. Esos números caben sobradamente en una MCU clase ARM Cortex-M con 256--384\,kB de Flash y 64--128\,kB de RAM, todavía más después de una cuantización a 8 bits. El LSTM-AE y el Transformer-AE quedan como alternativas razonables si se acepta ejecutar el detector en un concentrador (Raspberry Pi o equivalente) a cambio de un mejor $\aucpr$. LightGBM, aunque ligero, no es la recomendación final por la razón ya conocida: exige un corpus etiquetado de ataques que en este problema no existe.

\section{Resultados de la ampliación: técnicas de frontera}
\label{sec:resultados-ampliacion}

Esta sección recoge los resultados de los cinco experimentos de la ampliación (NB15--NB20), descritos en \cref{sec:diseno-ampliacion}. El hilo conductor es que el detector deja de evaluarse solo por su $\fone$ y pasa a medirse en cinco ejes adicionales: incertidumbre, generación, privacidad, robustez y eficiencia. Por honestidad metodológica, varios de estos experimentos arrojan \emph{resultados negativos o matizados}, que se reportan como tales por su valor científico.

\subsection{Incertidumbre: autoencoder variacional}
\label{sec:res-vae}

El $\beta$-VAE (\cref{sec:teo-vae}) alcanza $\fone=0{,}654$ y $\aucpr=0{,}456$, por debajo del Dense-AE determinista ($\fone=0{,}816$): es el precio del cuello de botella probabilístico y del muestreo Monte Carlo del \textit{score}. Su aportación no es la exactitud, sino dos propiedades nuevas. Primero, una \textbf{incertidumbre por alarma}: la \cref{fig:vae} muestra que la dispersión del \textit{score} sobre las muestras MC crece cerca del umbral ---justo en la zona de duda donde el operador debe intervenir---. Segundo, una \textbf{probabilidad calibrada}: tras un ajuste isotónico en validación, el error de calibración esperado en test es $\mathrm{ECE}=0{,}049$, de modo que ``probabilidad $0{,}8$'' significa de verdad un 80\,\% de aciertos. El modelo mantiene un \textit{footprint} \textit{edge} equivalente al Dense-AE (27,5\,kB, $0{,}0003$\,\si{\milli\second}/muestra).

\begin{figure}[H]
\centering
\includegraphics[width=0.49\textwidth]{vae_uncertainty.png}
\includegraphics[width=0.49\textwidth]{vae_reliability.png}
\caption{VAE probabilístico. Izquierda: la incertidumbre (desviación típica MC) crece cerca del umbral. Derecha: \emph{reliability diagram}; la cercanía a la diagonal indica buena calibración ($\mathrm{ECE}=0{,}049$).}
\label{fig:vae}
\end{figure}

\subsection{Modelos generativos: difusión}
\label{sec:res-diffusion}

El DDPM (\cref{sec:teo-ddpm}) entrenado sobre consumo limpio reproduce la distribución de la potencia con realismo \emph{moderado} (coeficiente de Bhattacharyya $0{,}81$; autocorrelación de lag-1 $0{,}74$ frente a $0{,}85$ real). Sin embargo, como \textbf{detector por \textit{denoising}} sobre la señal univariante de potencia es prácticamente aleatorio ($\aucroc=0{,}52$), a la par que un AE de ventana equivalente ($0{,}54$). Lejos de ser un fracaso, este resultado negativo es \emph{revelador}: sin las características físicas multivariantes ---en particular $\vires$, ausente al trabajar solo con la potencia--- la detección generativa pierde casi todo su poder. Refuerza así la tesis central del trabajo: el valor reside en la invariante física, no solo en la sofisticación del modelo generativo. El banco \emph{diffusion-FDI} (ataques sobre 1\,000 portadoras generadas) se entrega como artefacto complementario.

\subsection{Privacidad: aprendizaje federado y privacidad diferencial}
\label{sec:res-federated}

El \textbf{aprendizaje federado} entrena el Dense-AE sobre 6 clientes no-IID sin centralizar datos crudos y alcanza $\fone=0{,}812$ ($\aucpr=0{,}835$), \emph{igualando ---e incluso superando levemente---} a la línea base centralizada ($\fone=0{,}791$); el efecto regularizador de promediar modelos compensa la heterogeneidad. La \cref{fig:federated} (izquierda) muestra la convergencia de FedAvg ronda a ronda. Sobre esa base, \textbf{DP-SGD} introduce el compromiso privacidad--utilidad de la \cref{tab:dp}: al endurecer la garantía (menor $\varepsilon$), el $\aucpr$ cae de forma marcada ($\sim$0,52 frente a 0,84 sin privacidad), mientras que el $\fone$ se mantiene en torno a $0{,}72$--$0{,}74$. El presupuesto $\varepsilon$ se convierte así en un parámetro de cumplimiento que el operador fija según su apetito de riesgo, análogo a la recalibración de coste del \cref{sec:resultados-recalibracion}. Los valores de $\varepsilon$ reportados se obtienen con la cota RDP de orden entero del mecanismo gaussiano submuestreado (\cref{eq:rdp}); es una cota \emph{superior conservadora} ---contadores numéricos como los de \textit{TensorFlow Privacy} u \textit{Opacus} darían $\varepsilon$ algo menores---, por lo que la privacidad efectiva es \emph{al menos} la indicada.

\begin{table}[H]
\centering
\caption{Compromiso privacidad--utilidad de DP-SGD ($\delta=10^{-5}$, recorte $C=1$). A mayor multiplicador de ruido $\sigma$, menor $\varepsilon$ (más privacidad) y menor utilidad.}
\label{tab:dp}
\footnotesize
\begin{tabular}{cccc}
\toprule
$\sigma$ & $\varepsilon$ & $\fone$ & $\aucpr$ \\
\midrule
--- (sin DP)  & $\infty$  & 0{,}812 & 0{,}835 \\
0{,}6 & 15{,}95 & 0{,}721 & 0{,}526 \\
1{,}0 & 5{,}11  & 0{,}719 & 0{,}515 \\
1{,}5 & 2{,}44  & 0{,}719 & 0{,}482 \\
2{,}5 & 1{,}23  & 0{,}743 & 0{,}551 \\
\bottomrule
\end{tabular}
\end{table}

\begin{figure}[H]
\centering
\includegraphics[width=0.49\textwidth]{federated_convergence.png}
\includegraphics[width=0.49\textwidth]{dp_privacy_utility.png}
\caption{Izquierda: convergencia de FedAvg frente al centralizado. Derecha: frontera privacidad ($\varepsilon$, escala logarítmica) frente a utilidad bajo DP-SGD.}
\label{fig:federated}
\end{figure}

\subsection{Robustez adversaria de gradiente}
\label{sec:res-adversarial}

Sometido a ataques de caja blanca (\cref{sec:teo-adv}), el detector resulta \textbf{razonablemente robusto} en el espacio de características: a un presupuesto grande ($\epsilon=1{,}2$ en unidades de IQR), PGD solo evade el 13\,\% de los ataques previamente detectados y FGSM apenas el 1{,}5\,\% (\cref{fig:adv}, izquierda). La evasión depende de la severidad ---los ataques de baja severidad se ocultan más (17\,\%) que los de alta (11\,\%)---, lo que reencuentra por la vía del gradiente la frontera económica del \cref{sec:resultados-adversarial}.

Dos hallazgos sobre la defensa. Primero, un \emph{resultado negativo honesto}: el \textbf{adversarial training} de un AE no supervisado \emph{no} mejora la robustez a evasión (la empeora ligeramente, ganancia media $-0{,}05$), porque el modelo nunca ve ataques y solo suaviza su superficie de reconstrucción. Segundo, la defensa que \emph{sí} funciona es el \textbf{ensemble}: un ataque diseñado contra el detector A apenas transfiere a un detector B independiente (que aún detecta el 91\,\%), de modo que el conjunto A$\lor$B detecta el 96\,\% (\cref{fig:adv}, derecha). Esto valida cuantitativamente el \textit{stacking} del \cref{sec:resultados-ensemble} como mecanismo de robustez.

\begin{figure}[H]
\centering
\includegraphics[width=0.49\textwidth]{adversarial_robustness_gradient.png}
\includegraphics[width=0.49\textwidth]{adversarial_transfer_ensemble.png}
\caption{Robustez adversaria. Izquierda: tasa de detección frente al presupuesto $\epsilon$ (PGD $>$ FGSM en fuerza de ataque). Derecha: la evasión no transfiere ---el ensemble A$\lor$B detecta el 96\,\% del ataque diseñado contra A---.}
\label{fig:adv}
\end{figure}

\subsection{Eficiencia edge: cuantización y pruning}
\label{sec:res-edge-comp}

La \cref{tab:compresion} cuantifica la compresión del Dense-AE. Conviene leerla con cuidado: la pequeña caída de $\fone$ entre el modelo Keras ($0{,}817$) y su exportación a TFLite en \code{float32} ($0{,}784$) se debe a la \emph{conversión} (recalibración del umbral sobre las reconstrucciones del intérprete), no a la cuantización; tomando esa exportación como línea base, la cuantización \emph{float16} y \emph{dynamic-range} int8 es prácticamente \textbf{sin pérdida} ($0{,}784\to0{,}781$) mientras reduce el modelo a 14,1\,kB. El \emph{pruning} al 90\,\% de \textit{sparsity} lo deja en 4,4\,kB comprimido (gzip) con $\fone=0{,}74$ ---relevante para actualizaciones OTA a millones de contadores con ancho de banda limitado---. La latencia se mantiene en $\sim$$0{,}0016$\,\si{\milli\second}/muestra. El factor de compresión int8 es modesto ($\times 1{,}6$) precisamente porque el modelo ya nacía diminuto ($27$\,kB); el margen real lo aporta el \textit{pruning}.

\begin{table}[H]
\centering
\caption{Compresión del Dense-AE. PTQ vía \code{TFLiteConverter}; \emph{pruning} por magnitud con tamaño comprimido (gzip). $\fone$ recalibrado en validación para cada variante.}
\label{tab:compresion}
\footnotesize
\begin{tabular}{llcc}
\toprule
Método & Variante & Tamaño (kB) & $\fone$ \\
\midrule
Referencia & float32 (Keras)        & 27{,}0       & 0{,}817 \\
\midrule
TFLite     & float32 (línea base)   & 30{,}8       & 0{,}784 \\
PTQ        & float16                & 18{,}8       & 0{,}784 \\
PTQ        & dynamic int8           & \textbf{14{,}1} & 0{,}781 \\
PTQ        & full int8              & 17{,}2       & 0{,}741 \\
Pruning    & 70\,\% sparsity (gzip) & 10{,}0       & 0{,}741 \\
Pruning    & 90\,\% sparsity (gzip) & \textbf{4{,}4} & 0{,}742 \\
\bottomrule
\end{tabular}
\end{table}

\subsection{Síntesis multidimensional}
\label{sec:res-radar}

La \cref{fig:radar} resume el cambio de perspectiva: el mismo núcleo Dense-AE/VAE se evalúa ahora en cinco ejes. Ningún detector domina en todos ---el Dense-AE lidera en detección y eficiencia, el VAE aporta incertidumbre, el federado aporta privacidad y el ensemble aporta robustez---, lo que confirma que la elección de qué desplegar es multiobjetivo y no se reduce a un único número.

\begin{figure}[H]
\centering
\includegraphics[width=0.62\textwidth]{comparison_radar.png}
\caption{Perfil multidimensional de detectores representativos (detección, incertidumbre, privacidad, robustez y eficiencia). El valor de la ampliación es cubrir ejes que el $\fone$ puntual no captura.}
\label{fig:radar}
\end{figure}

\section{Optimización de hiperparámetros}
\label{sec:resultados-optimizacion}

Como cierre experimental se realizó una \textbf{búsqueda sistemática de
hiperparámetros} sobre los detectores, siguiendo un \emph{benchmark-optimization
loop}: (i) se fija el \emph{baseline} de cada modelo; (ii) se reproduce ese
\emph{baseline} dentro de un \emph{harness} de búsqueda autocontenido
---verificación previa imprescindible: el \emph{Isolation Forest} y el Dense-AE
se reprodujeron exactamente antes de buscar---; (iii) se exploran variantes por
muestreo aleatorio sobre la región de hiperparámetros, registrando cada una en un
CSV \emph{append-only}; y (iv) se promueve la mejor y se \textbf{re-valida a
escala completa} antes de confirmarla. El \emph{correctness gate} es el mismo
contrato de no-trampa de todo el trabajo: entrenamiento sobre \code{train\_clean},
calibración de umbral y de $(W,K)$ \emph{solo} en validación, evaluación en test,
sin fuga de etiquetas. Para acelerar la búsqueda, el entrenamiento y la
puntuación se realizan sobre submuestras; los finalistas se re-entrenan y evalúan
a escala completa, que es la cifra reportada en la \cref{tab:optim}.

\begin{table}[H]
\centering
\caption{Optimización de hiperparámetros: $\fone$ y $\aucpr$ antes (\emph{base})
y después (\emph{opt}) de la búsqueda, bajo el mismo contrato. Mejora neta en
negrita.}
\label{tab:optim}
\footnotesize
\begin{tabular}{lcccccc}
\toprule
Detector & $\fone$ base & $\fone$ opt & $\Delta\fone$ & $\aucpr$ base & $\aucpr$ opt & $\Delta\aucpr$ \\
\midrule
CNN-AE         & 0{,}435 & \textbf{0{,}643} & \textbf{+0{,}207} & 0{,}432 & 0{,}694 & +0{,}261 \\
Transformer-AE & 0{,}604 & \textbf{0{,}660} & +0{,}057 & 0{,}584 & 0{,}692 & +0{,}108 \\
VAE            & 0{,}693 & \textbf{0{,}747} & +0{,}054 & 0{,}500 & 0{,}758 & +0{,}258 \\
LSTM-AE        & 0{,}745 & 0{,}745$^{\dagger}$ & $\approx$0 & 0{,}881 & 0{,}881 & $\approx$0 \\
IF             & 0{,}754 & \textbf{0{,}795} & +0{,}041 & 0{,}782 & 0{,}853 & +0{,}070 \\
Dense-AE       & 0{,}816 & 0{,}815 & $\approx$0 & 0{,}855 & 0{,}855 & $\approx$0 \\
\bottomrule
\end{tabular}
\\[2pt]
{\footnotesize $^{\dagger}$ El mejor candidato del LSTM mejoraba sobre el subconjunto de búsqueda pero \emph{regresó} ($0{,}738$) al re-validar a escala completa; se conserva el \emph{baseline}. La re-validación detectó así un sobreajuste al subconjunto.}
\end{table}

Las mejoras provienen de tres palancas principales: (a) usar el conjunto de
características físicas \emph{completo} (17, incluyendo $\vires$ y derivadas) en
los autoencoders convolucional y \emph{transformer}, que en su versión original
empleaban un subconjunto de 14; (b) ponderar el error de reconstrucción solo por
las características \emph{discriminantes} ($\mathrm{AUC}>0{,}5$) en lugar de
incluir también las anti-discriminantes; y (c) en el \emph{Isolation Forest},
aumentar \code{max\_samples} de 256 a 2048, que mejora de forma marcada la
calidad de la puntuación de anomalía.

El hallazgo más llamativo es el \textbf{CNN-AE}, cuyo $\fone$ pasa de $0{,}435$ a
$0{,}643$ (un salto de $0{,}21$ puntos) y su $\aucpr$ de $0{,}43$ a $0{,}69$: el
\emph{baseline} estaba claramente infra-optimizado, y bastó corregir el conjunto
de características y la ponderación del \emph{score} para situarlo en línea con el
resto de autoencoders. El \textbf{VAE} mejora sobre todo en $\aucpr$ ($+0{,}26$)
al puntuar con la probabilidad de reconstrucción \emph{sin} ponderación espuria y
recalibrar el umbral. El \textbf{LSTM-AE}, en cambio, ilustra el valor de la re-validación: un candidato
que parecía mejorar sobre el subconjunto resultó un sobreajuste y se descartó al
evaluarlo a escala completa, conservándose el \emph{baseline} ya afinado. En el extremo opuesto,
el \textbf{Dense-AE} ---ya muy afinado en el trabajo original--- no admite mejora
apreciable, lo que es en sí mismo un resultado: confirma que era una elección de
diseño acertada. En conjunto, la optimización \textbf{no altera las conclusiones
cualitativas} del trabajo (se mantiene el orden relativo y la brecha
supervisado/no-supervisado), pero estrecha el margen de los detectores más
débiles.

\begin{observacion}[Reproducibilidad]
Toda la búsqueda es reproducible mediante los \emph{scripts} del directorio
\code{optimization/} (\code{opt\_pointwise.py}, \code{opt\_window.py},
\code{finalize.py}, \code{consolidate.py}), que reutilizan el módulo común
\code{tfg\_common.py} ---y por tanto el mismo contrato de datos y protocolo de
evaluación que los \emph{notebooks}---. Los \emph{notebooks} de partida se
conservan intactos como \emph{baseline} documentado; la optimización es una capa
adicional, no una reescritura.
\end{observacion}

